\documentclass[12pt]{article}
\usepackage[utf8]{inputenc}
\usepackage[T1]{fontenc}
\usepackage{amsmath,amssymb}
\usepackage{geometry}
\geometry{margin=2.5cm}

\begin{document}

\noindent $N=\{e,\sigma,\sigma^2\} \leq S_3$.

\bigskip
\noindent $H=\{e,\tau\}$ \quad $\tau = \begin{pmatrix}1&2&3\\2&1&3\end{pmatrix}$, \ $\tau^2=e$.

\noindent $\Rightarrow H\leq S_3$.

\bigskip
\noindent $N \trianglelefteq S_3$ \quad dar \quad $H\not\trianglelefteq S_3$.

\bigskip
\noindent $a\in S_3$ şi $x\in H \overset{?}{\Rightarrow} axa^{-1}\in H$.\footnote{în original apare ,,$x\in N$'', dar calculul care urmează testează normalitatea lui $H$ (folosind $\tau\in H$) şi obține un rezultat care nu e în $H$ -- corectat aici la $H$ pentru consecvență cu restul argumentului.}

\[
\begin{pmatrix}1&2&3\\2&3&1\end{pmatrix}\begin{pmatrix}1&2&3\\2&1&3\end{pmatrix}\begin{pmatrix}1&2&3\\3&1&2\end{pmatrix} = \begin{pmatrix}1&2&3\\1&3&2\end{pmatrix} \notin H.
\]

\bigskip
\noindent\underline{Exp} \ 1) Fie $(G,\cdot,e)$=grup.

\noindent $\{e\}\trianglelefteq G$. \quad şi \quad $G\trianglelefteq G$.

\bigskip
\noindent (2).\ $GL_2(\mathbb{R}) = \{A\in M_2(\mathbb{R}) \mid \det(A)\neq 0\}$.

\noindent $SL_2(\mathbb{R}) = \{A\in M_2(\mathbb{R}) \mid \det(A)=1\}$.

\noindent $(M_2(\mathbb{R}),\cdot)$ = monoid.

\noindent $GL_2(\mathbb{R})$ = p.s.\ a lui $M_2(\mathbb{R})$ în raport cu $\cdot$

\noindent $(GL_2(\mathbb{R}),\cdot)$ = grup necomutativ.

\noindent $SL_2(\mathbb{R}) \trianglelefteq GL_2(\mathbb{R})$.

\bigskip
\noindent N.\ Becheanu, C.\ Vraciu -- P.S.\ de teoria grupurilor.\ Tip.\ Univ.\footnote{referință bibliografică; abrevierea ,,P.S.'' e neclară -- ar putea fi ,,Probleme [Selectate/de Seminar] de teoria grupurilor''.}

\bigskip
\noindent\underline{Lema}: Fie $(G,\cdot)$=grup.\ şi $R\subseteq G\times G$ rel.\ de ech.\ pe $G$.

\noindent U.A.S.E.

\begin{enumerate}
\item[(1)] $R$ compatibilă cu $\cdot$. \ $\big(xRx'$ şi $yRy' \Rightarrow xyRx'y'\big)$.
\item[(2)] $R$ e compatibilă la st.\ şi la dr.\ cu $\cdot$

\noindent \big($xRy \Rightarrow zxRzy$ şi $xzRyz$, $(\forall)\,z\in G$\big).
\end{enumerate}

\bigskip
\noindent\underline{Dem} $(1)\Rightarrow(2)$.\ $zRz \overset{(1)}{\Longrightarrow} (xzRyz)\wedge(zxRzy)$.

\bigskip
\noindent $(2) \Rightarrow (1)$. \quad $xRx' \Rightarrow xy'Rx'y'$

\noindent $yRy' \Rightarrow xyRxy'$

\noindent $\Rightarrow xyRx'y'$.

\end{document}
